% 原章第59—63页：附录5.B与全部9道习题；原第64页为出版商标志页。
\section*{附录 5.B\quad 分块 HSIC 的渐近零分布}
\phantomsection\addcontentsline{toc}{section}{附录 5.B：分块 HSIC 的渐近零分布}
\markright{附录 5.B}

注意，
\begin{equation}
\Var(\mathrm{BHSIC})=\frac Sn\Var(\widehat c_b).
\tag{5.B.1}
\end{equation}
由于 $\mathrm{var}(S\widehat c_b)\approx\mathrm{var}(W)$，且 $W$ 的分布为
$\sum_{i=1}^{\infty}\sum_{j=1}^{\infty}\lambda_i\eta_j(N_{i,j}^2-1)$，得到
\begin{equation}
\Var(\widehat c_b)\approx\frac1{S^2}\mathrm{var}(W).
\tag{5.B.2}
\end{equation}
众所周知，$\chi_{(1)}^2$ 的方差为 2。由于 $N_{i,j}$ 独立同分布于 $N(0,1)$，得到
\begin{equation}
\Var(W)=2\sum_{i=1}^{\infty}\sum_{j=1}^{\infty}\lambda_i^2\eta_j^2.
\tag{5.B.3}
\end{equation}
令 $\widehat K_x(X,X')$ 和 $\widehat K_y(Y,Y')$ 为中心化核矩阵。于是
\begin{equation}
\tr(\widehat K_x^2)\to\sum_{i=1}^{\infty}\lambda_i^2,
\qquad\text{且}\qquad
\tr(\widehat K_y^2)\to\sum_{j=1}^{\infty}\eta_j^2.
\tag{5.B.4}
\end{equation}
回顾以下关系：
\begin{equation}
\begin{aligned}
\tr(\widehat K_x^2)
&=\sum_{i=1}^S\sum_{j=1}^Sk_x^2(x_i,x_j)
=S^2\frac1{S^2}\sum_{i=1}^S\sum_{j=1}^Sk_x^2(x_i,x_j)\\
&\to S^2E_{XX'}[\widehat k_x^2(X,X')],\\
\text{以及}\qquad
\tr(\widehat K_y^2)&\to S^2E_{YY'}[\widehat k_y^2(Y,Y')].
\end{aligned}
\tag{5.B.5}
\end{equation}
将式 (5.B.4) 和 (5.B.5) 代入式 (5.B.3)，得到
\begin{equation}
\mathrm{var}(W)=2S^4E_{XX'}[\widehat k_x^2(X,X')]E_{YY'}[\widehat k_y^2(Y,Y')].
\tag{5.B.6}
\end{equation}
结合式 (5.B.1) 和 (5.B.2)，得到
\begin{equation}
\Var(\mathrm{BHSIC})=\frac1{ns}\Var(W).
\tag{5.B.7}
\end{equation}
因此，定义
\begin{equation}
T_{BI}=\sqrt{ns}\frac{\mathrm{BHSIC}}{\Var(W)}.
\tag{5.B.8}
\end{equation}
在相互独立的零假设下，$T_{BI}$ 渐近服从 $N(0,1)$ 分布。

\section*{习题}
\phantomsection\addcontentsline{toc}{section}{习题}
\markright{习题}
\begin{enumerate}[label=\textbf{习题 \arabic*.},leftmargin=4.7em]
\item 假定 $k$ 为线性核，$P$ 为正态分布 $N(0,1)$，$Q$ 为二项分布。计算最大均值差异 $\mathrm{MMD}_k(P,Q)$。

\item 令 $A$ 为 $n\times n$ 维矩阵，$B$ 为 $m\times m$ 维矩阵。证明
\[
|A\otimes B|=|A|^n|B|^m,\qquad
\Tr(A\otimes B)=\Tr(A)\Tr(B),
\]
其中，$\Tr$ 表示矩阵的迹。

\item 证明
\[
\langle A,B\rangle_{\mathrm{HS}}
=\sum_{i\in I}\sum_{j\in J}\langle Av_j,e_i\rangle_F\langle Bv_j,e_i\rangle_F.
\]

\item 证明
\[
\sum_{j=1}^n\sum_{l=1}^n\sum_{i\ne(j,l)}^nk_x(x_j,x_i)k_y(y_i,y_l)
=1_n^T\widehat K_x\widehat K_y1_n,
\]
以及
\[
\sum_{j=1}^n\sum_{i\ne j}k_x(x_j,x_i)k_y(y_i,y_j)
=\tr(\widehat K_x\widehat K_y).
\]

\item 证明
\[
\begin{aligned}
&\sum_{(i,q)\in\mathrm i_2^n}k_x(x_i,x_q)\sum_{(j,l)\in\mathrm i_2^n}k_y(y_j,y_l)\\
&=\sum_{(i,q,j)\in\mathrm i_3^n}k_x(x_i,x_q)k_y(y_j,y_i)\\
&\quad+\sum_{(i,q,l)\in\mathrm i_3^n}k_x(x_i,x_q)k_y(y_q,y_l)\\
&\quad+\sum_{(i,j,l)\in\mathrm i_3^n}k_x(x_i,x_j)k_y(y_j,y_l)\\
&\quad+\sum_{(q,j,l)\in\mathrm i_3^n}k_x(x_l,x_q)k_y(y_j,y_l)\\
&\quad+\sum_{(i,j,q,l)\in\mathrm i_4^n}k_x(x_i,x_q)k_y(y_j,y_l)\\
&\quad+\sum_{(i,q)\in\mathrm i_2^n}k_x(x_i,x_q)k_y(y_i,y_q)\\
&\quad+\sum_{(i,q)\in\mathrm i_2^n}k_x(x_i,x_q)k_y(y_i,y_q).
\end{aligned}
\]

\item 证明正文中的式 (5.92)—(5.94)：
\begin{equation}
\begin{aligned}
\widehat\Delta_i
={}&\left[W_i^TX_0(\sigma_i^2X_0^TX_0)^{-1}X_0^TW_i\right]^{-1}
W_i^TX_0(\sigma_i^2X_0^TX_0)^{-1}X_0^Ty_i\\
={}&\left[W_i^TX_0(X_0^TX_0)^{-1}X_0^TW_i\right]^{-1}
W_i^TX_0(X_0^TX_0)^{-1}X_0^Ty_i,
\end{aligned}
\tag{5.92}
\end{equation}
以及
\begin{equation}
\Lambda_i=\mathrm{var}(\widehat\Delta_i)
=N\widehat\sigma_i^2\left[W_i^TX_0(X_0^TX_0)^{-1}X_0^TW_i\right]^{-1},
\tag{5.93}
\end{equation}
其中，$\widehat\sigma_i^2$ 估计为
\begin{equation}
\widehat\sigma_i^2=\frac1{NJ-m_iJ+1-k_i}
(y_i-W_i\widehat\Delta_i)^T(y_i-W_i\widehat\Delta_i),
\tag{5.94}
\end{equation}
$m_i$ 与 $k_i$ 分别为第 $i$ 个方程中所含的内生变量和外生变量数目。

\item 证明正文中问题 (5.99) 的最优解为
\[
\begin{aligned}
\widehat\gamma_i&=\left[A_1-\sum_{j=1}^JB_1(j)(B_2(j)+\rho I)^{-1}A_2(j)\right]^{-1}D,\\
D&=D_1-\sum_{j=1}^JB_1(j)(B_2(j)+\rho I)^{-1}D_2(j),\\
\eta_i^{(j)}&=[B_2(j)+\rho I]^{-1}[D_2(j)-A_2(j)\widehat\gamma_i],
\end{aligned}
\]
其中
\[
\begin{aligned}
A_1&=\sum_{j=1}^J(Y_{-i}^{(j)})^TX(\sigma_{ii}^{(j)}X^TX)^{-1}X^TY_{-i}^{(j)},\\
A_2(j)&=X^TX(\sigma_{ii}^{(j)}X^TX)^{-1}X^TY_{-i}^{(j)},\\
B_1(j)&=(Y_{-i}^{(j)})^TX(\sigma_{ii}^{(j)}X^TX)^{-1}X^TX,\\
B_2(j)&=X^TX(\sigma_{ii}^{(j)}X^TX)^{-1}X^TX,\\
D_1&=\sum_{j=1}^J(Y_{-i}^{(j)})^TX(\sigma_{ii}^{(j)}X^TX)^{-1}X^Ty_i^{(j)}
+\rho(Z_{i\gamma}^{(l)}-\mu_\gamma^{(l)}),\\
D_2(j)&=X^TX(\sigma_{ii}^{(j)}X^TX)^{-1}X^Ty_i^{(j)}
+\rho(Z_i^{(j)(l)}-\mu_j^{(l)}).
\end{aligned}
\]

\item 证明问题 (5.101) 的最优解为
\[
Z_{\lambda i}^{(l+1)}=\sgn(\Delta_{\gamma i}^{(l+1)}+\mu_\gamma^{(l)})
\left(|\Delta_{\gamma i}^{(l+1)}+\mu_\gamma^{(l)}|-\frac\lambda\rho\right)_+,
\]
其中
\[
|x|_+=\begin{cases}x,&x\geq0,\\0,&x<0.\end{cases}
\]

\item 证明式 (5.105) 推出（习题 8）
\[
\left(1+\frac{\lambda\sqrt J}{\rho\lVert Z_{ki}\rVert_2}\right)\lVert Z_{ki}\rVert_2
=\lVert\Delta_{ki}^{(l+1)}-\mu_{ki}^{(l)}\rVert_2.
\]
\end{enumerate}
